% !TEX program = pdflatex
\documentclass[aspectratio=169]{beamer}

\usetheme{Madrid}
\usecolortheme{default}

\usepackage{amsmath,amssymb,mathtools}
\usepackage{physics}
\usepackage{bm}

\title{Derivative of the Variational Energy}
\subtitle{Deriving the Monte Carlo Gradient Formula}
\author{Morten Hjorth-Jensen}
\date{Spring 2026}

\newcommand{\avg}[1]{\left\langle #1 \right\rangle}
\newcommand{\psia}{\psi_{\bm{\alpha}}}
\newcommand{\dalpha}{\partial_{\alpha}}
\newcommand{\El}{E_L}

\begin{document}

%========================================================
\begin{frame}
\titlepage
\end{frame}

%========================================================
\begin{frame}{Setup and notation}
Let $\psia(x)$ be a (real-valued) variational wavefunction depending on parameters
$\bm{\alpha}=(\alpha_1,\dots,\alpha_P)$ and $x$ denote the many-body coordinates (continuous or discrete).

\begin{block}{Variational energy}
\[
E(\bm{\alpha})=\frac{\bra{\psia}H\ket{\psia}}{\bra{\psia}\psia\rangle}
=
\frac{\int dx\, \psia(x)\, [H\psia](x)}{\int dx\, |\psia(x)|^2}.
\]
\end{block}

Define the normalization:
\[
\mathcal{N}(\bm{\alpha}) := \int dx\, |\psia(x)|^2.
\]
\end{frame}

%========================================================
\begin{frame}{Local energy and sampling distribution}
Introduce the \textbf{local energy}
\[
\El(x;\bm{\alpha}) := \frac{[H\psia](x)}{\psia(x)}.
\]

Define the normalized probability density
\[
p_{\bm{\alpha}}(x) := \frac{|\psia(x)|^2}{\mathcal{N}(\bm{\alpha})},
\qquad
\int dx\, p_{\bm{\alpha}}(x)=1.
\]

Then the energy can be written as an expectation value:
\[
E(\bm{\alpha}) = \int dx\, p_{\bm{\alpha}}(x)\, \El(x;\bm{\alpha})
= \avg{\El}.
\]

\begin{block}{Monte Carlo viewpoint}
We estimate $E(\bm{\alpha})$ by sampling $x\sim p_{\bm{\alpha}}$ and averaging $\El(x)$.
\end{block}
\end{frame}

%========================================================
\section{Derivative of the energy}

%========================================================
\begin{frame}{Differentiate $E(\bm{\alpha})$ in quotient form}
Start from
\[
E(\bm{\alpha})=\frac{\int dx\, \psia(x)\, [H\psia](x)}{\int dx\, \psia(x)^2}
\]
(assume $\psia$ real; complex case is similar with conjugation).

Let
\[
A(\bm{\alpha}) := \int dx\, \psia\, H\psia,
\qquad
B(\bm{\alpha}) := \int dx\, \psia^2.
\]
Then $E=A/B$ and
\[
\dalpha E = \frac{(\dalpha A)B - A(\dalpha B)}{B^2}.
\]
\end{frame}

%========================================================
\begin{frame}{Compute $\partial_\alpha A$ and $\partial_\alpha B$}
Compute
\[
\dalpha B = \dalpha \int dx\, \psia^2
= \int dx\, 2\psia\, (\dalpha \psia).
\]

For $A$:
\[
\dalpha A
= \dalpha \int dx\, \psia\, (H\psia)
= \int dx\, \Big[(\dalpha \psia)(H\psia) + \psia\, H(\dalpha \psia)\Big].
\]

\begin{block}{Use Hermiticity of $H$}
Assuming suitable boundary conditions, $H$ is Hermitian:
\[
\int dx\, \psia\, H(\dalpha \psia) = \int dx\, (\dalpha \psia)\, H\psia.
\]
Thus
\[
\dalpha A = 2\int dx\, (\dalpha \psia)\, H\psia.
\]
\end{block}
\end{frame}

%========================================================
\begin{frame}{Insert into quotient rule}
Insert $\dalpha A$ and $\dalpha B$ into
\[
\dalpha E = \frac{(\dalpha A)B - A(\dalpha B)}{B^2}.
\]

Using $A = EB$:
\[
\dalpha E
= \frac{2\left(\int dx\, (\dalpha \psia)\, H\psia\right)B
- (EB)\left(\int dx\, 2\psia\, \dalpha\psia\right)}{B^2}.
\]

Cancel $B$:
\[
\dalpha E
= 2\left[
\frac{\int dx\, (\dalpha \psia)\, H\psia}{\int dx\, \psia^2}
- E\,\frac{\int dx\, \psia\, \dalpha\psia}{\int dx\, \psia^2}
\right].
\]
\end{frame}

%========================================================
\begin{frame}{Rewrite using $p_{\alpha}(x)$ and local energy}
Divide and multiply by $\psia$:
\[
\frac{\int dx\, (\dalpha \psia)\, H\psia}{\int dx\, \psia^2}
=
\frac{\int dx\, \psia^2 \left(\frac{\dalpha \psia}{\psia}\right)
\left(\frac{H\psia}{\psia}\right)}{\int dx\, \psia^2}.
\]

Recognize:
\[
\frac{H\psia}{\psia}=\El,
\qquad
\frac{\dalpha\psia}{\psia} = \dalpha \ln \psia.
\]

Thus:
\[
\dalpha E
= 2\left[
\avg{\left(\frac{\dalpha\psia}{\psia}\right)\El}
- E\,\avg{\frac{\dalpha\psia}{\psia}}
\right],
\qquad E=\avg{\El}.
\]
\end{frame}

%========================================================
\begin{frame}{Final covariance form}
Since $E=\avg{\El}$,
\[
\dalpha E
=
2\left(
\avg{\left(\frac{\dalpha\psia}{\psia}\right)\El}
-\avg{\frac{\dalpha\psia}{\psia}}\avg{\El}
\right).
\]

\begin{block}{Compact form (covariance)}
Define the ``log-derivative observable''
\[
O_\alpha(x):=\frac{\dalpha\psia(x)}{\psia(x)}=\dalpha \ln \psia(x),
\]
then
\[
\boxed{
\dalpha E = 2\,\mathrm{Cov}\!\left(O_\alpha,\El\right)
=2\left(\avg{O_\alpha \El}-\avg{O_\alpha}\avg{\El}\right).
}
\]
\end{block}
\end{frame}

%========================================================
\section{Interpretation and the ``three integrals''}

%========================================================
\begin{frame}{The three expectation values / integrals}
Written explicitly in terms of $p_{\bm{\alpha}}(x)$, the gradient contains:

\[
\avg{O_\alpha \El}=\int dx\, p_{\bm{\alpha}}(x)\,
\left(\frac{\dalpha\psia(x)}{\psia(x)}\right)\El(x),
\]

\[
\avg{O_\alpha}=\int dx\, p_{\bm{\alpha}}(x)\,
\left(\frac{\dalpha\psia(x)}{\psia(x)}\right),
\]

\[
\avg{\El}=\int dx\, p_{\bm{\alpha}}(x)\,\El(x).
\]

\begin{block}{Why this is useful}
Each integral is an expectation value with respect to the \emph{same} sampling distribution
$p_{\bm{\alpha}}(x)\propto \psia(x)^2$, so all three are estimated from the \emph{same} Monte Carlo run.
\end{block}
\end{frame}

%========================================================
\begin{frame}{Practical Monte Carlo estimator}
Given samples $\{x^{(m)}\}_{m=1}^M \sim p_{\bm{\alpha}}$, define
\[
O_\alpha^{(m)} := \left.\frac{\dalpha\psia}{\psia}\right|_{x^{(m)}},
\qquad
\El^{(m)} := \El(x^{(m)}).
\]

Then the unbiased (up to sampling error) estimator is
\[
\widehat{\dalpha E}
=
2\left(
\frac{1}{M}\sum_{m=1}^M O_\alpha^{(m)}\El^{(m)}
-\left(\frac{1}{M}\sum_{m=1}^M O_\alpha^{(m)}\right)
\left(\frac{1}{M}\sum_{m=1}^M \El^{(m)}\right)
\right).
\]

\begin{block}{Interpretation}
The gradient is proportional to the covariance between the log-derivative of the wavefunction
and the local energy.
\end{block}
\end{frame}

%========================================================
\begin{frame}{Remarks and extensions}
\begin{itemize}
\item If $\psia$ is complex, replace $\psia^2$ by $|\psia|^2$ and
\[
O_\alpha \equiv \partial_\alpha \ln \psia \quad \text{is generally complex.}
\]
The physically relevant gradient involves the \emph{real part}:
\[
\partial_\alpha E = 2\,\Re\left(\avg{O_\alpha^\ast \El}-\avg{O_\alpha^\ast}\avg{\El}\right).
\]
\item The same derivation applies to excited-state functionals and generalized cost functions.
\item This identity underlies stochastic reconfiguration / natural gradient methods in VMC.
\end{itemize}
\end{frame}

%========================================================
\begin{frame}{Summary}
\begin{itemize}
\item Start from the Rayleigh quotient $E=\langle \psi|H|\psi\rangle/\langle\psi|\psi\rangle$.
\item Differentiate using quotient rule and Hermiticity of $H$.
\item Introduce $p(x)\propto |\psi(x)|^2$ and the local energy $\El=(H\psi)/\psi$.
\item Obtain the VMC gradient:
\[
\partial_\alpha E =
2\left(\avg{O_\alpha \El}-\avg{O_\alpha}\avg{\El}\right).
\]
\end{itemize}
\end{frame}

\end{document}
